\documentclass[10pt]{article}
\usepackage[margin=1in]{geometry}
\usepackage[T1]{fontenc}
\usepackage{lmodern,microtype,amsmath,amssymb,amsthm}
\theoremstyle{definition}
\newtheorem{assumption}{Assumption}
\begin{document}
\section*{Heavy-Tailed Two-Point Zeroth-Order Optimization}
We consider minimizing a differentiable, possibly nonconvex function $f:\mathbb{R}^d\to\mathbb{R}$, where $d\geq1$, with a lower bound $\ell\in\mathbb{R}$ and given starting point $x^0$ satisfying $f(x^0)-\ell\leq\Delta$, where $\Delta\geq0$. The oracle returns two scalar values $F(x,\xi)$ and $F(y,\xi)$ using a shared sample $\xi\sim\nu$, with independent samples across calls. Each sample function $F(\cdot,\xi)$ has an $L$-Lipschitz gradient, where $L>0$; these gradients are unbiased for $\nabla f$ and have centered $p$th moments bounded by $\sigma^p$, where $1<p\leq2$ and $\sigma\geq0$. For $\varepsilon>0$, the goal is a measurable output $\bar{x}$ satisfying $\mathbb{E}\|\nabla f(\bar{x})\|\leq\varepsilon$, with expectation over sampling and independent algorithm randomness.

The optimization problem is
\[
\inf_{x\in\mathbb{R}^d} f(x),
\]
where $\|\cdot\|$ denotes the Euclidean norm. The algorithm receives $d$, the measurable sample space $\Omega$, $L$, $\Delta$, $\sigma$, $p$, $\varepsilon$, and $x^0\in\mathbb{R}^d$.

Our assumptions:
\begin{assumption}[Objective and initial gap]
The function $f$ is differentiable everywhere. There exists $\ell\in\mathbb{R}$ such that
\[
\begin{aligned}
f(x)&\geq\ell &&\text{for every }x\in\mathbb{R}^d,\\
f(x^0)-\ell&\leq\Delta.
\end{aligned}
\]
\end{assumption}

\begin{assumption}[Sample smoothness]
There exist a probability measure $\nu$ on $\Omega$ and jointly measurable maps
\[
F:\mathbb{R}^d\times\Omega\to\mathbb{R},
\qquad
G:\mathbb{R}^d\times\Omega\to\mathbb{R}^d.
\]
For every $\omega\in\Omega$, the function $F(\cdot,\omega)$ is differentiable everywhere, and for every $x,y\in\mathbb{R}^d$,
\[
\begin{aligned}
\nabla_x F(x,\omega)&=G(x,\omega),\\
\|G(x,\omega)-G(y,\omega)\|&\leq L\|x-y\|.
\end{aligned}
\]
\end{assumption}

\begin{assumption}[Unbiased gradients and heavy-tailed noise]
For every $x\in\mathbb{R}^d$, the map $G(x,\cdot)$ is Bochner integrable under $\nu$, and $\|G(x,\cdot)-\nabla f(x)\|^p$ is integrable under $\nu$. They satisfy
\[
\begin{aligned}
\mathbb{E}_{\xi\sim\nu}[G(x,\xi)]&=\nabla f(x),\\
\mathbb{E}_{\xi\sim\nu}\bigl[\|G(x,\xi)-\nabla f(x)\|^p\bigr]&\leq\sigma^p.
\end{aligned}
\]
\end{assumption}

At call $k\geq0$, the algorithm queries any pair $(x_k,y_k)\in\mathbb{R}^d\times\mathbb{R}^d$ and receives
\[
\bigl(F(x_k,\xi_k),F(y_k,\xi_k)\bigr).
\]
The sample variables $(\xi_k)_{k\geq0}$ are measurable, mutually independent, and each has law $\nu$. They are independent of the algorithm's randomness. Thus both values in a reply use the same sample, while distinct calls use independent samples. One oracle call returns this pair of scalar values; the budget counts such paired calls.

After the allotted calls, the output $\bar{x}\in\mathbb{R}^d$ must be measurable, its gradient norm must be integrable, and it must satisfy
\[
\mathbb{E}_{P}\bigl[\|\nabla f(\bar{x})\|\bigr]\leq\varepsilon,
\]
where $P$ is the joint law of the samples and the independent algorithm randomness.
\end{document}
